\chapter*{Справочные таблицы}\addcontentsline{toc}{chapter}{Справочные таблицы}\markright{Таблицы}

\section*{Таблица умножения}

\begingroup\renewcommand{\arraystretch}{1.25}\setlength{\tabcolsep}{6pt}\arrayrulecolor{Teal}

\begin{center}\begin{tabular}{|c|r|r|r|r|r|r|r|r|r|}\hline

\(\times\)&1&2&3&4&5&6&7&8&9\\\hline

1&1&2&3&4&5&6&7&8&9\\\hline

2&2&4&6&8&10&12&14&16&18\\\hline

3&3&6&9&12&15&18&21&24&27\\\hline

4&4&8&12&16&20&24&28&32&36\\\hline

5&5&10&15&20&25&30&35&40&45\\\hline

6&6&12&18&24&30&36&42&48&54\\\hline

7&7&14&21&28&35&42&49&56&63\\\hline

8&8&16&24&32&40&48&56&64&72\\\hline

9&9&18&27&36&45&54&63&72&81\\\hline

\end{tabular}\end{center}\endgroup

\par\noindent\begin{minipage}{\linewidth}\section*{Тригонометрические значения (радианы)}

\begingroup\small\setlength{\tabcolsep}{4pt}\arrayrulecolor{Teal}\begin{center}\begin{tabular}{|c|c|c|c|c|c|c|c|c|}\hline

\tablemath{x}&\tablemath{0}&\tablemath{\frac\pi6}&\tablemath{\frac\pi4}&\tablemath{\frac\pi3}&\tablemath{\frac\pi2}&\tablemath{\pi}&\tablemath{\frac{3\pi}2}&\tablemath{2\pi}\\\hline

\tablemath{\sin x}&\tablemath{0}&\tablemath{\frac12}&\tablemath{\frac{\sqrt2}2}&\tablemath{\frac{\sqrt3}2}&\tablemath{1}&\tablemath{0}&\tablemath{-1}&\tablemath{0}\\\hline

\tablemath{\cos x}&\tablemath{1}&\tablemath{\frac{\sqrt3}2}&\tablemath{\frac{\sqrt2}2}&\tablemath{\frac12}&\tablemath{0}&\tablemath{-1}&\tablemath{0}&\tablemath{1}\\\hline

\tablemath{\operatorname{tg}x}&\tablemath{0}&\tablemath{\frac{\sqrt3}3}&\tablemath{1}&\tablemath{\sqrt3}&---&\tablemath{0}&---&\tablemath{0}\\\hline

\tablemath{\operatorname{ctg}x}&---&\tablemath{\sqrt3}&\tablemath{1}&\tablemath{\frac{\sqrt3}3}&\tablemath{0}&---&\tablemath{0}&---\\\hline

\end{tabular}\end{center}\endgroup\end{minipage}\par

\clearpage\section*{Производные базовых функций}

\begingroup\small\arrayrulecolor{Teal}\setlength{\tabcolsep}{5pt}\begin{longtable}{|p{.24\linewidth}|p{.28\linewidth}|p{.38\linewidth}|}\hline\tablemath{f(x)}&\tablemath{f'(x)}&Условия\\\hline\endfirsthead\hline\tablemath{f(x)}&\tablemath{f'(x)}&Условия\\\hline\endhead

\tablemath{c}&\tablemath{0}&\(c\) — постоянная; \(x\in\mathbb R\).\\\hline

\tablemath{x}&\tablemath{1}&\(x\in\mathbb R\).\\\hline

\tablemath{x^n}&\tablemath{nx^{n-1}}&\(n\in\mathbb N\), \(n\ge2\); \(x\in\mathbb R\).\\\hline

\tablemath{x^\alpha}&\tablemath{\alpha x^{\alpha-1}}&\(\alpha\in\mathbb R\) постоянна; \(x>0\).\\\hline

\tablemath{\sqrt x}&\tablemath{\frac1{2\sqrt x}}&\(x>0\).\\\hline

\tablemath{\frac1x}&\tablemath{-\frac1{x^2}}&\(x\ne0\).\\\hline

\tablemath{e^x}&\tablemath{e^x}&\(x\in\mathbb R\).\\\hline

\tablemath{a^x}&\tablemath{a^x\ln a}&\(a>0\) постоянно; \(x\in\mathbb R\).\\\hline

\tablemath{\ln x}&\tablemath{\frac1x}&\(x>0\).\\\hline

\tablemath{\log_a x}&\tablemath{\frac1{x\ln a}}&\(a>0\), \(a\ne1\), \(x>0\).\\\hline

\tablemath{\sin x}&\tablemath{\cos x}&\(x\in\mathbb R\), радианы.\\\hline

\tablemath{\cos x}&\tablemath{-\sin x}&\(x\in\mathbb R\), радианы.\\\hline

\tablemath{\operatorname{tg}x}&\tablemath{\frac1{\cos^2x}}&\(\cos x\ne0\), радианы.\\\hline

\tablemath{\operatorname{ctg}x}&\tablemath{-\frac1{\sin^2x}}&\(\sin x\ne0\), радианы.\\\hline

\tablemath{\arcsin x}&\tablemath{\frac1{\sqrt{1-x^2}}}&\(-1<x<1\).\\\hline

\tablemath{\arccos x}&\tablemath{-\frac1{\sqrt{1-x^2}}}&\(-1<x<1\).\\\hline

\tablemath{\operatorname{arctg}x}&\tablemath{\frac1{1+x^2}}&\(x\in\mathbb R\).\\\hline

\tablemath{\operatorname{sh}x}&\tablemath{\operatorname{ch}x}&\(x\in\mathbb R\).\\\hline

\tablemath{\operatorname{ch}x}&\tablemath{\operatorname{sh}x}&\(x\in\mathbb R\).\\\hline

\tablemath{\operatorname{th}x}&\tablemath{\frac1{\operatorname{ch}^2x}}&\(x\in\mathbb R\).\\\hline

\tablemath{\operatorname{cth}x}&\tablemath{-\frac1{\operatorname{sh}^2x}}&\(x\ne0\).\\\hline

\end{longtable}\endgroup
